% vim: tw=78 ai spell spelllang=cs
\documentclass{article}

%\usepackage{a4wide}
\usepackage{graphicx}
\usepackage[utf8x]{inputenc}
\usepackage[czech]{babel}
\usepackage[IL2]{fontenc}

\pagestyle{empty}

\begin{document}
\begin{flushright}
\today
\end{flushright}
\begin{center}
\Large\textbf{Posudek oponenta na bakalářskou práci \\
\normalsize%
\textit{Vít Šesták:} Trasování hodnot proměnných v JVM}
\end{center}

Dle zadání měla bakalářská práce implementovat nástroj, který bude schopen u
příkazů Java bytekódu říci, jaká může být hodnota jednotlivých výrazů a odkud
se sem dostala. Text práce je na 27 stranách strukturovaný do 4 částí, které v
součtu obsahují 18 podkapitol. V úvodu autor krátce popisuje stav a motivaci
pro práci. Další kapitola se snaží popisovat ,,Teoretické základy'', které
mají být využity v další kapitole ,,Implementace''. Práci ukončuje závěr
obsahující i popis práce do budoucna. \textbf{Rozsah} uvedených kapitol je z
pohledu bakalářské práce i zadání dostatečný.

\textbf{Obsahově} přesto některé pasáže práce působí jako záznam volných
myšlenkových pochodů bez návaznosti na sebe. Rozdělení do tolika podkapitol
toto ještě zhoršuje -- práce je nepřehledná a bez vysvětlení použitých pojmů.
Popř. jsou-li některé vysvětleny/odkázány, tak až po jejich použití (kap.
1.4.8 popisuje pojem \emph{invokedynamic}, použit je ale už v 1.4.2). Sekce
,,Pojmy'' na str. 3 nemělo být třeba. Vše mělo být součástí textu v nějaké
logické posloupnosti.  Práce dle mého názoru také nedostatečně uvádí do
problematiky (např. na začátku není uvedeno, proč se instrumentuje, přičemž se
okolo v textu píše o instrumentaci) a nedostatečně popisuje teoretické
základy, tedy kapitola 1 je až příliš obsahově strohá. Autor zde mohl na mnoha
místech použít grafického znázornění, např. sekce 1.7.2 k tomu přímo vybízí. A
když už autor sáhne k příkladu, neuvede ho kompletní: v sekci 2.10.2 by bylo
nanejvýše vhodné uvést, s jakými parametry byl program spuštěn a na jakém
zdrojovém kódu, abychom dostali uvedený výstup.  Také je toho pro autora až
příliš ,,zřejmé'', pro čtenáře až tak ne (např. sekce 1.3, odst. 2, věta 2.).
Autor se měl více věnovat sběru podkladů pro svá tvrzení, aby se v textu tak
často neobjevovala slova ,,pravděpodobně'' a ,,nejspíš'' bez uvedení faktů, na
základě kterých se to autor domnívá (např. kap.  3, odst. 3).  Literatura
obsahuje mnoho zvláštních referencí do jednoho webu (ke každé instrukci
zmíněné v práci) -- toto mohl být jeden odkaz. Referovat sebe sama jako
,,autora'' záznamu v bugzille vypadá trochu podivně.

Po \textbf{jazykové} stránce práce obsahuje malé množství překlepů (např.
,,budou považován'', ,,je bude'' oboje v 1.1, ,,přistp'' v 1.3.2). Větší počet
prohřešků je na poli typografie. Práce má zbytečně velké mezery mezi odstavci,
přetoky řádků, např. v 1.7.25 a 2.8.2 (zde zůstala část textu někde na
tiskovém válci). Text je nepatrně širší než obvyklých 80 znaků na řádek.
Tabulka na str. 30 je zvláštně vysázená.

\textbf{Praktická část} práce je v přiloženém \texttt{tar.gz} souboru.
Přestože je v textu minimalistická dokumentace použití nástroje, při dodržení
pokynů nefunguje. Třída k vypisování historie je totiž
\texttt{com.v6ak.adFontes.library.Debug}, nikoliv
\texttt{com.v6ak.adFontes.Debug} (viz sekce 2.10).  Kód je dobře čitelný a
jeví se být funkční.

\begin{itemize}
\item 2.9: uvádí podporu nástroje jen zastaralé Javy 6 a nižší. Proč?
\item 3.1.3: v práci (opět) chybí konkrétní údaje či lépe: grafy. Přímo se
naskýtá otázka jak velké ono zpomalení je? Bylo to měřeno?
\item 3.1.4: jsou zmíněny pluginy. Podporuje je implementace?
\item K 3.1, další možné rozšíření: program je nyní postavený tak, že běží
dynamicky a vypisuje jen hodnotu jednoho běhu programu. Bylo by možné přístup
rozšířit tak, aby fungoval staticky (analýza bez JVM)? Jak obtížné by to bylo?
\end{itemize}

Na závěr konstatuji, že práce \textbf{splňuje zadání a požadavky na
bakalářskou práci}. Navrhuji práci \textbf{uznat} jako bakalářskou, a vzhledem
k výše uvedeným okolnostem doporučuji hodnocení stupněm C--D v závislosti na
obhajobě.

\vspace{3em}
\begin{flushright}
Jiří Slabý
\end{flushright}

\end{document}
